\unitchapter{1}{4}{Communication}{p1-04}

\unitmeta{Expected questions: \textasciitilde5 (10 marks) · Target: 5/5}

\begin{sylbox}

\begin{itemize}
\tightlist
\item[$\square$]
  Communication: meaning, types and characteristics
\item[$\square$]
  Effective communication: verbal \& non-verbal, inter-cultural \& group communication, classroom communication
\item[$\square$]
  Barriers to effective communication
\item[$\square$]
  Mass-media and society
\end{itemize}

\end{sylbox}

\needspace{191pt}

\hypertarget{p1-04--1-meaning}{%
\section{Meaning}\label{p1-04--1-meaning}}

At its simplest, communication succeeds when the receiver understands what the sender meant. Every model
in this unit is an attempt to explain how that shared meaning is achieved --- and why it so often is not.
Latin \emph{communis / communicare} = "to make common / to share". Communication is the process of transmitting
information, ideas and feelings to create \textbf{shared understanding}.

\needspace{125pt}

\hypertarget{p1-04--2-models-of-communication}{%
\section{Models of Communication}\label{p1-04--2-models-of-communication}}

Models simplify a complex process so that its parts can be studied. Early \emph{linear} models (Aristotle,
Lasswell, Shannon--Weaver) treat communication as a one-way transmission. \emph{Interactive} models add feedback,
and \emph{transactional} models see both parties sending and receiving at once. The single most examined idea is
\textbf{noise}, introduced by Shannon and Weaver: anything that distorts the message on its way.

\begin{figure}[!htbp]
\centering
\begin{tikzpicture}[node distance=17mm]
  \node[slate, minimum width=24mm] (S) {\textbf{Sender}\\\footnotesize source};
  \node[grey, minimum width=22mm, right=of S] (M) {\textbf{Message}};
  \node[sand, minimum width=22mm, right=of M] (C) {\textbf{Channel}};
  \node[sage, minimum width=24mm, right=of C] (R) {\textbf{Receiver}\\\footnotesize destination};
  \draw[arr] (S) -- node[lbl, above] {encoding} (M);
  \draw[arr] (M) -- (C);
  \draw[arr] (C) -- node[lbl, above] {decoding} (R);
  \draw[dasharr] (R.south) -- ++(0,-9mm) -| node[lbl, below, pos=0.25] {feedback} (S.south);
  \node[cloudnode, above=8mm of C, minimum width=22mm, minimum height=11mm, fill=frose] (N) {\footnotesize\textbf{Noise}};
  \draw[decorate, decoration={zigzag, segment length=4pt, amplitude=1.5pt}, fwine, line width=0.6pt, ->] (N) -- (C);
\end{tikzpicture}
\caption{The general model of communication. Shannon and Weaver's contribution was the noise source, which
can distort the message at any point on the channel.}
\end{figure}


\begingroup\small

\par\medskip\noindent\hfil\begin{tabular}{@{}
  >{\raggedright\arraybackslash}p{(\columnwidth - 4\tabcolsep) * \real{0.2333}}
  >{\raggedright\arraybackslash}p{(\columnwidth - 4\tabcolsep) * \real{0.2879}}
  >{\raggedright\arraybackslash}p{(\columnwidth - 4\tabcolsep) * \real{0.4787}}@{}}
\toprule\noalign{}
\begin{minipage}[b]{\linewidth}\raggedright
\textbf{Model}
\end{minipage} & \begin{minipage}[b]{\linewidth}\raggedright
\textbf{Type}
\end{minipage} & \begin{minipage}[b]{\linewidth}\raggedright
\textbf{Key point}
\end{minipage} \\
\midrule\noalign{}
\textbf{Aristotle} & Linear, oldest & Speaker → Speech → Audience (persuasion: ethos, pathos, logos) \\
\textbf{Lasswell (1948)} & Linear & \emph{Who says What in which Channel to Whom with what Effect} \\
\textbf{Shannon--Weaver (1949)} & Linear, "mother of all models" & Introduced \textbf{noise}; information source, transmitter, channel, receiver, destination \\
\textbf{Berlo SMCR (1960)} & Linear & Source, Message, Channel, Receiver \\
\textbf{Osgood--Schramm} & Circular & Both are encoder-interpreter-decoder \\
\textbf{Schramm (1954)} & Interactive & \textbf{Field of experience} must overlap \\
\textbf{Westley--MacLean} & Mass comm & Gatekeeper role \\
\textbf{Dance (helical)} & Transactional & Communication grows like a helix \\
\textbf{Barnlund} & Transactional & Simultaneous sending \& receiving \\
\bottomrule\noalign{}
\end{tabular}\hfil\null\par\medskip


\endgroup

\begin{figure}[!htbp]
\centering
\begin{adjustbox}{max width=\linewidth}
\begin{tikzpicture}
  \begin{scope}[blend group=multiply]
    \fill[fslate] (0,0) circle (2.3);
    \fill[fsand] (3.0,0) circle (2.3);
  \end{scope}
  \draw[fgink, line width=0.5pt] (0,0) circle (2.3) (3.0,0) circle (2.3);
  \node[align=center] at (-1.0,0.4) {\textbf{Sender}\\\footnotesize encoder};
  \node[align=center] at (4.0,0.4) {\textbf{Receiver}\\\footnotesize decoder};
  \node[align=center, font=\footnotesize\itshape] at (1.5,0) {shared\\meaning};
  \node[lbl] at (-1.0,-1.25) {sender's field\\of experience};
  \node[lbl] at (4.0,-1.25) {receiver's field\\of experience};
  \draw[arr] (-0.6,2.65) to[bend left=22] node[lbl, above] {message} (3.6,2.65);
  \draw[dasharr] (3.6,-2.65) to[bend left=22] node[lbl, below] {feedback} (-0.6,-2.65);
\end{tikzpicture}
\end{adjustbox}
\caption{Schramm's model: a message can be understood only to the extent that the fields of experience of
sender and receiver overlap.}
\end{figure}


\needspace{125pt}

\hypertarget{p1-04--3-types-of-communication}{%
\section{Types of Communication}\label{p1-04--3-types-of-communication}}

Communication can be classified by the channel used, the number of people involved, the direction of flow
in an organisation, and its degree of formality. A single act --- a principal\textquotesingle s circular to teachers, for
example --- is written, downward, formal and group communication all at once.

\begin{figure}[!htbp]
\centering
\begin{adjustbox}{max width=\linewidth}
\begin{tikzpicture}
  \foreach \x/\col/\head/\items in {
      0/fslate/{By channel}/{verbal — oral, written\\non-verbal},
      1/fsage/{By participants}/{intrapersonal (self)\\interpersonal (two)\\group, public\\mass (through media)},
      2/fsand/{By direction}/{upward, downward\\horizontal (lateral)\\diagonal\\grapevine (informal)},
      3/frose/{By formality}/{formal\\informal}} {
    \node[box, fill=\col, minimum width=34mm] at (\x*3.75,0) {\textbf{\head}};
    \node[plain, text width=32mm, align=left, font=\footnotesize, anchor=north, minimum height=17mm] at (\x*3.75,-0.6) {\items};
  }
\end{tikzpicture}
\end{adjustbox}
\caption{Four independent ways of classifying communication.}
\end{figure}


\needspace{125pt}

\hypertarget{p1-04--non-verbal-communication-frequently-asked}{%
\subsection{Non-verbal communication (frequently asked)}\label{p1-04--non-verbal-communication-frequently-asked}}

Much of what we communicate is carried not by words but by the body, the voice, space and time. The
technical names (kinesics, proxemics, haptics, chronemics) are frequently asked; Edward T. Hall\textquotesingle s zones of
personal space are illustrated below.

\begin{figure}[!htbp]
\centering
\begin{tikzpicture}
  \foreach \r/\c in {3.4/fgrey, 2.5/fslate, 1.6/fsage, 0.8/frose} \filldraw[fill=\c, draw=fgink, line width=0.45pt] (0,0) circle (\r);
  \fill[fgink] (0,0.12) circle (0.12); \draw[fgink, line width=1.2pt] (0,0) -- (0,-0.35);
  \foreach \r/\name/\d in {0.8/Intimate/{0–1.5 ft}, 1.6/Personal/{1.5–4 ft}, 2.5/Social/{4–12 ft}, 3.4/Public/{12 ft +}} {
    \draw[thinline] (30:\r) -- (5.0,{\r*0.5+0.1}) node[right, align=left, font=\small] {\textbf{\name}\ \ \footnotesize\d};
  }
\end{tikzpicture}
\caption{Edward T. Hall's proxemic zones: the distance people keep communicates the closeness of the
relationship.}
\end{figure}


\begingroup\small

\par\medskip\noindent\hfil\begin{tabular}{@{}
  >{\raggedright\arraybackslash}p{(\columnwidth - 2\tabcolsep) * \real{0.2211}}
  >{\raggedright\arraybackslash}p{(\columnwidth - 2\tabcolsep) * \real{0.4767}}@{}}
\toprule\noalign{}
\begin{minipage}[b]{\linewidth}\raggedright
\textbf{Kind}
\end{minipage} & \begin{minipage}[b]{\linewidth}\raggedright
\textbf{Meaning}
\end{minipage} \\
\midrule\noalign{}
\textbf{Kinesics} & Body movements, gestures, facial expressions, posture \\
\textbf{Proxemics} & Use of space / distance (Edward T. Hall) \\
\textbf{Haptics} & Touch \\
\textbf{Chronemics} & Use of time \\
\textbf{Paralanguage / Vocalics} & Tone, pitch, volume, pauses --- \emph{how} it is said \\
\textbf{Oculesics} & Eye contact \\
\textbf{Artifactics} & Clothing, appearance, objects \\
\textbf{Olfactics} & Smell \\
\bottomrule\noalign{}
\end{tabular}\hfil\null\par\medskip


\endgroup

Hall\textquotesingle s proxemic zones: Intimate (0--1.5 ft), Personal (1.5--4 ft), Social (4--12 ft), Public (12 ft+).

\textbf{Mehrabian\textquotesingle s 7-38-55 rule} (for feelings/attitudes): 7\% words, 38\% tone of voice, 55\% body language.

\needspace{95pt}

\hypertarget{p1-04--classroom-communication}{%
\subsection{Classroom Communication}\label{p1-04--classroom-communication}}

\begin{itemize}
\tightlist
\item
  Mostly \textbf{interpersonal + group}; should be two-way, with feedback.
\item
  \textbf{Flanders\textquotesingle{} Interaction Analysis Category System (FIACS)} --- 10 categories (7 teacher talk, 2 student talk, 1 silence/confusion); teacher talk = indirect (1--4) + direct (5--7).
\item
  Effective classroom communication: clarity, eye contact, questioning, active listening, empathy, using examples.
\end{itemize}

\needspace{125pt}

\hypertarget{p1-04--4-barriers-to-communication}{%
\section{Barriers to Communication}\label{p1-04--4-barriers-to-communication}}

A barrier is anything that prevents the intended meaning from reaching the receiver. Questions usually
describe a situation --- jargon in a lecture, a noisy hall, a student\textquotesingle s prejudice --- and ask you to name the
barrier, so learn the categories with an example of each.

\begin{figure}[!htbp]
\centering
\begin{tikzpicture}
  \node[circ, minimum size=26mm, fill=fwine!15, font=\small\bfseries] (B) at (0,0) {Barriers};
  \foreach \a/\h/\t/\col [count=\k] in {90/Physical/{noise, distance,\\poor light}/fslate, 38.6/Semantic/{jargon,\\ambiguity}/fsage,
      -12.9/Psychological/{prejudice,\\emotion}/fsand, -64.3/Cultural/{values, norms,\\gestures}/frose,
      -115.7/Organisational/{hierarchy,\\rigid rules}/fslate, -167.1/Technical/{faulty\\equipment}/fsage, 141.4/Perceptual/{stereotyping,\\selectivity}/fsand} {
    \node[box, fill=\col, text width=24mm, font=\footnotesize] (n\k) at (\a:3.6) {\textbf{\h}\\\t};
    \draw[thinline] (B) -- (n\k);
  }
\end{tikzpicture}
\caption{Seven families of barriers to effective communication.}
\end{figure}


7 Cs of effective communication: \textbf{Clear, Concise, Concrete, Correct, Coherent, Complete, Courteous}.

\needspace{125pt}

\hypertarget{p1-04--5-inter-cultural--group-communication}{%
\section{Inter-cultural \& Group Communication}\label{p1-04--5-inter-cultural--group-communication}}

When sender and receiver come from different cultures, their fields of experience overlap less, and
misunderstanding becomes more likely. Groups add their own dynamics: they pass through recognisable stages
before they work well together.

\begin{itemize}
\tightlist
\item
  \textbf{Hofstede\textquotesingle s cultural dimensions}: power distance, individualism--collectivism, masculinity--femininity, uncertainty avoidance, long-term orientation, indulgence.
\item
  \textbf{High-context} cultures (India, Japan) rely on implicit cues; \textbf{low-context} (USA, Germany) are explicit (E.T. Hall).
\item
  Ethnocentrism = judging other cultures by one\textquotesingle s own --- barrier.
\item
  Tuckman\textquotesingle s group stages: \textbf{Forming → Storming → Norming → Performing → Adjourning}.
\end{itemize}

\begin{figure}[!htbp]
\centering
\begin{tikzpicture}
  \foreach \i/\n/\t/\col in {0/Forming/{polite,\\uncertain}/fslate, 1/Storming/{conflict over\\roles}/frose, 2/Norming/{rules and\\cohesion}/fsage,
      3/Performing/{productive\\teamwork}/fsand, 4/Adjourning/{task ends,\\group disbands}/fgrey} {
    \node[box, fill=\col, text width=22mm, font=\footnotesize, anchor=south west] (t\i) at (\i*2.75,\i*0.55) {\textbf{\n}\\\t};
  }
  \foreach \i [evaluate=\i as \j using int(\i+1)] in {0,...,3} \draw[arr] (t\i.east) -- (t\j.west);
\end{tikzpicture}
\caption{Tuckman's stages of group development.}
\end{figure}


\needspace{125pt}

\hypertarget{p1-04--6-mass-media--society}{%
\section{Mass Media \& Society}\label{p1-04--6-mass-media--society}}

Theories of mass communication have swung between two views of the audience: as passive targets of
powerful media (the "magic bullet") and as active users who select, interpret and resist. The two-step flow
model introduced the influential middle layer of \emph{opinion leaders}.

\begin{figure}[!htbp]
\centering
\begin{tikzpicture}
  \node[slate, minimum width=26mm, minimum height=12mm] (M) at (0,0) {\textbf{Mass media}};
  \foreach \y/\n in {1.3/O1, -1.3/O2} \node[circ, fill=fsand, minimum size=14mm, font=\footnotesize, align=center] (\n) at (4,\y) {opinion\\leader};
  \foreach \y [count=\k] in {2.1,1.3,0.5} \node[circ, minimum size=6mm] (a\k) at (7,\y) {};
  \foreach \y [count=\k] in {-0.5,-1.3,-2.1} \node[circ, minimum size=6mm] (b\k) at (7,\y) {};
  \draw[arr] (M) -- node[lbl, above, sloped] {step 1} (O1); \draw[arr] (M) -- (O2);
  \foreach \k in {1,2,3} \draw[arr] (O1) -- (a\k);
  \foreach \k in {1,2,3} \draw[arr] (O2) -- (b\k);
  \node[lbl] at (5.7,2.65) {step 2}; \node[capt] at (7,-2.85) {the public};
\end{tikzpicture}
\caption{Lazarsfeld and Katz's two-step flow: media influence passes through opinion leaders to the wider public.}
\end{figure}


\begingroup\small

\par\medskip\noindent\hfil\begin{tabular}{@{}
  >{\raggedright\arraybackslash}p{(\columnwidth - 2\tabcolsep) * \real{0.4311}}
  >{\raggedright\arraybackslash}p{(\columnwidth - 2\tabcolsep) * \real{0.5411}}@{}}
\toprule\noalign{}
\begin{minipage}[b]{\linewidth}\raggedright
\textbf{Theory}
\end{minipage} & \begin{minipage}[b]{\linewidth}\raggedright
\textbf{Idea}
\end{minipage} \\
\midrule\noalign{}
\textbf{Hypodermic needle / Magic bullet} & Media has direct, powerful effect on passive audience \\
\textbf{Two-step flow} (Lazarsfeld \& Katz) & Media → opinion leaders → public \\
\textbf{Agenda-setting} (McCombs \& Shaw) & Media tells us \emph{what to think about} \\
\textbf{Uses \& Gratifications} (Katz, Blumler) & Active audience uses media to satisfy needs \\
\textbf{Cultivation} (George Gerbner) & Long TV exposure shapes perception of reality ("mean world") \\
\textbf{Spiral of silence} (Noelle-Neumann) & Minority opinion holders stay silent \\
\textbf{Medium is the message} (Marshall McLuhan) & "Global village"; hot vs cool media \\
\textbf{Gatekeeping} (Kurt Lewin, David Manning White) & Editors filter news \\
\bottomrule\noalign{}
\end{tabular}\hfil\null\par\medskip


\endgroup

\textbf{Functions of mass media (Lasswell + Wright)}: surveillance, correlation, transmission of culture, entertainment.
Indian landmarks: First newspaper --- \emph{Hicky\textquotesingle s Bengal Gazette} (1780, James Augustus Hicky); Doordarshan (1959);
All India Radio (1936, named; Akashvani 1957); Press Council of India (1966); Prasar Bharati (1997).

\needspace{226pt}

\hypertarget{p1-04--7-deeper-dive--listening-feedback-networks--classroom-talk}{%
\section{Deeper Dive --- Listening, Feedback, Networks \& Classroom Talk}\label{p1-04--7-deeper-dive--listening-feedback-networks--classroom-talk}}

\needspace{190pt}

\hypertarget{p1-04--71-listening}{%
\subsection{Listening}\label{p1-04--71-listening}}

Listening is the neglected half of communication. A message is complete only when it has been heard,
understood and responded to, which is why the stages of listening end with \emph{responding}.

Hearing is physiological; \textbf{listening} is psychological (attending, understanding, remembering, responding).

\begingroup\small

\par\medskip\noindent\hfil\begin{tabular}{@{}
  >{\raggedright\arraybackslash}p{(\columnwidth - 2\tabcolsep) * \real{0.2767}}
  >{\raggedright\arraybackslash}p{(\columnwidth - 2\tabcolsep) * \real{0.6494}}@{}}
\toprule\noalign{}
\begin{minipage}[b]{\linewidth}\raggedright
\textbf{Type of listening}
\end{minipage} & \begin{minipage}[b]{\linewidth}\raggedright
\textbf{Purpose}
\end{minipage} \\
\midrule\noalign{}
Active / empathetic & Understand the speaker\textquotesingle s feelings; paraphrase, nod, ask clarifying questions \\
Critical / evaluative & Judge the message, logic and evidence \\
Informational / comprehensive & Learn and retain content (lectures) \\
Appreciative & Enjoyment (music, poetry) \\
Selective & Hearing only what one wants --- a barrier \\
\bottomrule\noalign{}
\end{tabular}\hfil\null\par\medskip


\endgroup

Stages (HURIER model, Brownell): \textbf{Hearing → Understanding → Remembering → Interpreting → Evaluating → Responding}.

\needspace{125pt}

\hypertarget{p1-04--72-feedback}{%
\subsection{Feedback}\label{p1-04--72-feedback}}

Feedback turns a monologue into a dialogue. For the sender it confirms whether the message was understood;
for the receiver it is information about how others see them --- the idea captured by the Johari window.

\begin{itemize}
\tightlist
\item
  \textbf{Positive} (reinforces), \textbf{negative / corrective} (points out errors), \textbf{descriptive} (specific, non-judgemental), \textbf{evaluative} (judgement).
\item
  Good feedback is specific, timely, focused on behaviour (not the person), and actionable.
\item
  \textbf{Johari window} (Joseph Luft \& Harrington Ingham): self-awareness in communication.
\end{itemize}

\begin{figure}[!htbp]
\centering
\begin{adjustbox}{max width=\linewidth}
\begin{tikzpicture}
  \fill[fsage] (0,0) rectangle (3.6,2.2); \fill[fsand] (3.6,0) rectangle (7.2,2.2);
  \fill[fslate] (0,-2.2) rectangle (3.6,0); \fill[fgrey] (3.6,-2.2) rectangle (7.2,0);
  \draw[fgink, line width=0.55pt] (0,-2.2) rectangle (7.2,2.2) (3.6,-2.2) -- (3.6,2.2) (0,0) -- (7.2,0);
  \node[align=center] at (1.8,1.1) {\textbf{Open}\\\footnotesize arena};
  \node[align=center] at (5.4,1.1) {\textbf{Blind}\\\footnotesize spot};
  \node[align=center] at (1.8,-1.1) {\textbf{Hidden}\\\footnotesize façade};
  \node[align=center] at (5.4,-1.1) {\textbf{Unknown}};
  \node[capt] at (1.8,2.55) {known to self}; \node[capt] at (5.4,2.55) {not known to self};
  \node[capt, rotate=90] at (-0.4,1.1) {known to others}; \node[capt, rotate=90] at (-0.4,-1.1) {not known to others};
  \draw[arr, fwine] (4.6,1.9) -- node[lbl, above, text=fwine] {feedback} (3.75,1.9);
  \draw[arr, fwine] (1.8,-0.35) -- node[lbl, right, text=fwine] {self-disclosure} (1.8,0.25);
\end{tikzpicture}
\end{adjustbox}
\caption{The Johari window. Feedback from others shrinks the blind area; self-disclosure shrinks the hidden
area; both enlarge the open area in which communication is easiest.}
\end{figure}


\needspace{125pt}

\hypertarget{p1-04--73-communication-networks-small-groups}{%
\subsection{Communication Networks (small groups)}\label{p1-04--73-communication-networks-small-groups}}

Small-group experiments (Bavelas, Leavitt) showed that the \emph{pattern} of who may talk to whom affects speed,
accuracy, leadership and morale. Centralised networks are efficient for simple tasks; decentralised ones
suit complex problems and keep members more satisfied.

\begin{figure}[!htbp]
\centering
\begin{tikzpicture}[v/.style={circ, minimum size=6mm, inner sep=0pt, font=\scriptsize}, e/.style={thinline}]
  % wheel
  \begin{scope}[shift={(0,0)}]
    \node[v, fill=fwine!25] (c) at (0,0) {L};
    \foreach \a in {45,135,225,315} { \node[v] (w\a) at (\a:1.05) {}; \draw[e] (c) -- (w\a); }
    \node[capt] at (0,-1.6) {wheel};
  \end{scope}
  % chain
  \begin{scope}[shift={(3.2,0)}]
    \foreach \i in {0,...,4} \node[v] (k\i) at (0,{1.0-\i*0.5}) {};
    \foreach \i [evaluate=\i as \j using int(\i+1)] in {0,...,3} \draw[e] (k\i) -- (k\j);
    \node[capt] at (0,-1.6) {chain};
  \end{scope}
  % Y
  \begin{scope}[shift={(5.6,0)}]
    \node[v] (y1) at (-0.5,1.0) {}; \node[v] (y2) at (0.5,1.0) {}; \node[v] (y3) at (0,0.3) {};
    \node[v] (y4) at (0,-0.35) {}; \node[v] (y5) at (0,-1.0) {};
    \draw[e] (y1) -- (y3) -- (y2); \draw[e] (y3) -- (y4) -- (y5);
    \node[capt] at (0,-1.6) {Y};
  \end{scope}
  % circle
  \begin{scope}[shift={(8.2,0)}]
    \foreach \a in {90,162,234,306,18} \node[v] (r\a) at (\a:1.0) {};
    \draw[e] (r90) -- (r162) -- (r234) -- (r306) -- (r18) -- (r90);
    \node[capt] at (0,-1.6) {circle};
  \end{scope}
  % all channel
  \begin{scope}[shift={(11.2,0)}]
    \foreach \a in {90,162,234,306,18} \node[v] (q\a) at (\a:1.0) {};
    \foreach \a in {90,162,234,306,18} \foreach \b in {90,162,234,306,18} {\ifnum\a<\b \draw[e] (q\a) -- (q\b);\fi}
    \node[capt] at (0,-1.6) {all-channel};
  \end{scope}
  \draw[arrd] (-0.9,-2.25) -- node[lbl, below] {centralised\quad$\longleftrightarrow$\quad decentralised} (12.1,-2.25);
\end{tikzpicture}
\caption{Small-group communication networks. Centralised patterns (left) are fast and produce a clear
leader; decentralised ones (right) solve complex problems better and satisfy members more.}
\end{figure}


\begingroup\small

\par\medskip\noindent\hfil\begin{tabular}{@{}
  >{\raggedright\arraybackslash}p{(\columnwidth - 8\tabcolsep) * \real{0.1284}}
  >{\raggedright\arraybackslash}p{(\columnwidth - 8\tabcolsep) * \real{0.2425}}
  >{\raggedright\arraybackslash}p{(\columnwidth - 8\tabcolsep) * \real{0.1125}}
  >{\raggedright\arraybackslash}p{(\columnwidth - 8\tabcolsep) * \real{0.2287}}
  >{\raggedright\arraybackslash}p{(\columnwidth - 8\tabcolsep) * \real{0.2038}}@{}}
\toprule\noalign{}
\begin{minipage}[b]{\linewidth}\raggedright
\textbf{Network}
\end{minipage} & \begin{minipage}[b]{\linewidth}\raggedright
\textbf{Speed (simple tasks)}
\end{minipage} & \begin{minipage}[b]{\linewidth}\raggedright
\textbf{Accuracy}
\end{minipage} & \begin{minipage}[b]{\linewidth}\raggedright
\textbf{Member satisfaction}
\end{minipage} & \begin{minipage}[b]{\linewidth}\raggedright
\textbf{Leader emergence}
\end{minipage} \\
\midrule\noalign{}
Wheel & High & High & Low (except centre) & High \\
Chain & Moderate & High & Low & Moderate \\
Circle & Slow & Low & High & None \\
All-channel & Fast for complex tasks & Moderate & \textbf{Highest} & None \\
\bottomrule\noalign{}
\end{tabular}\hfil\null\par\medskip


\endgroup

\needspace{95pt}

\hypertarget{p1-04--74-effective-classroom-communication}{%
\subsection{Effective Classroom Communication}\label{p1-04--74-effective-classroom-communication}}

\begin{itemize}
\tightlist
\item
  Use \textbf{simple language}, examples from learners\textquotesingle{} lives, and checks for understanding.
\item
  Combine verbal with \textbf{non-verbal} cues (eye contact, gestures, movement).
\item
  \textbf{Two-way}: questions, discussions, feedback.
\item
  \textbf{Barriers} in class: noise, overcrowding, teacher\textquotesingle s monotone, unclear objectives, learner\textquotesingle s anxiety.
\item
  \textbf{Teacher immediacy}: behaviours that reduce psychological distance (smiling, using names) → better learning.
\end{itemize}

\needspace{95pt}

\hypertarget{p1-04--75-communication-ethics--digital-communication}{%
\subsection{Communication Ethics \& Digital Communication}\label{p1-04--75-communication-ethics--digital-communication}}

\begin{itemize}
\tightlist
\item
  Accuracy, honesty, respect for privacy, avoidance of plagiarism and misinformation.
\item
  \textbf{Fake news / misinformation} (false but shared without intent to deceive) vs \textbf{disinformation} (deliberately false).
\item
  \textbf{Netiquette}: rules of good online behaviour; ALL CAPS = shouting.
\item
  \textbf{Digital divide}: unequal access to ICT.
\end{itemize}

\needspace{95pt}

\hypertarget{p1-04--previous-year-questions-pyq-pattern}{%
\section{Previous Year Questions (PYQ pattern)}\label{p1-04--previous-year-questions-pyq-pattern}}

\begin{enumerate}
\def\labelenumi{\arabic{enumi}.}
\tightlist
\item
  "Who says what in which channel to whom with what effect" was given by: \textbf{Harold Lasswell}
\item
  Which model first introduced the concept of "noise"? \textbf{Shannon--Weaver}
\item
  The study of body movements as communication: \textbf{Kinesics}
\item
  Use of space in communication: \textbf{Proxemics}
\item
  The study of touch: \textbf{Haptics}
\item
  Tone and pitch of voice belong to: \textbf{Paralanguage}
\item
  Communication from subordinates to superiors: \textbf{Upward}
\item
  Informal communication network in an organisation: \textbf{Grapevine}
\item
  Jargon and ambiguous words cause which barrier? \textbf{Semantic}
\item
  Prejudice is a: \textbf{Psychological barrier}
\item
  "The medium is the message" --- \textbf{Marshall McLuhan}
\item
  Media determines "what to think about" --- \textbf{Agenda-setting theory}
\item
  Theory that the audience actively selects media to satisfy needs: \textbf{Uses and Gratifications}
\item
  Heavy TV viewers perceive the world as more dangerous --- \textbf{Cultivation theory (Gerbner)}
\item
  "Global village" --- \textbf{McLuhan}
\item
  Feedback in communication helps to: \textbf{verify that the message was understood as intended}
\item
  Classroom communication is basically: \textbf{Interactive / two-way, interpersonal + group}
\item
  FIACS was developed by: \textbf{Ned Flanders} (10 categories)
\item
  Arrange Tuckman\textquotesingle s stages: \textbf{Forming, Storming, Norming, Performing, Adjourning}
\item
  The first newspaper in India: \textbf{Hicky\textquotesingle s Bengal Gazette (1780)}
\item
  Communication with oneself: \textbf{Intrapersonal}
\item
  Statement: "Effective communication requires the sender and receiver to share a common field of experience." --- \textbf{True (Schramm)}
\end{enumerate}

\textbf{More practice questions}

\begin{enumerate}
\def\labelenumi{\arabic{enumi}.}
\setcounter{enumi}{22}
\tightlist
\item
  The first stage of the listening process: \textbf{Hearing (receiving)}
\item
  In the Johari window, information known to others but not to self is the: \textbf{blind area}
\item
  Self-disclosure reduces the: \textbf{hidden area}
\item
  The most centralised small-group network: \textbf{wheel}
\item
  Network with highest member satisfaction: \textbf{all-channel}
\item
  Feedback should be: \textbf{specific, timely and descriptive}
\item
  Deliberately false information spread to deceive: \textbf{disinformation}
\item
  Typing in ALL CAPS in online communication is considered: \textbf{shouting (poor netiquette)}
\item
  Teacher behaviours that reduce psychological distance: \textbf{immediacy behaviours}
\item
  Which is NOT a characteristic of effective classroom communication? (a) two-way (b) jargon-heavy (c) clear (d) learner-centred --- \textbf{Ans: (b)}
\end{enumerate}

\begin{revbox}

\begin{itemize}
\tightlist
\item
  Lasswell = 5 Ws · Shannon-Weaver = noise · Schramm = field of experience · Berlo = SMCR
\item
  Kinesics--body, Proxemics--space, Haptics--touch, Chronemics--time
\item
  Agenda setting (what to think about) vs Cultivation (TV shapes reality) vs Two-step (opinion leaders)
\end{itemize}

\end{revbox}
